% !TeX root = ./main.tex

\chapter{相关工作}\label{ch:related}

本章从检索增强生成、图增强检索、多模态文档理解和视觉语言模型四个方面梳理相关研究。

\section{检索增强生成（RAG）}

Lewis 等人在 NeurIPS 2020 提出的 RAG 框架证明了``参数知识 + 非参数知识''的组合方式能显著提升知识密集型 NLP 任务的准确性和可追溯性~\cite{lewis2020rag}。典型 RAG 流程将文档切分为文本块并存入 FAISS 等向量数据库~\cite{johnson2019faiss}，查询时通过语义相似度检索相关片段交由生成模型作答。然而标准 RAG 将文档处理为独立文本片段，难以处理跨段落、跨页面、多实体关联和证据链推理的复杂问题；在图文混排场景中，PDF 文本抽取或 OCR 阶段还易丢失版面信息、图像语义和表格结构。模块化 RAG 和自反思 RAG 等改进方案主要在文本检索框架内优化，未从根本上解决文档结构关系和多模态信息融合问题。

\section{图增强检索与 GraphRAG}

GraphRAG 引入图结构来组织内容，使检索能沿图关系扩展证据。Edge 等人提出的 GraphRAG 框架将文档构造成实体知识图谱并生成社区摘要，同时支持围绕实体的局部搜索（Local Search）和基于社区总结的全局搜索（Global Search）~\cite{edge2024graphrag}。微软官方实现进一步提供了基于 Leiden 算法的层级社区检测和 map-reduce 式全局回答流程~\cite{microsoft2025graphrag}。Bu 等人在 ACL 2025 提出了 Query-Driven Multimodal GraphRAG，强调根据查询语义动态构建局部图以提升跨模态推理~\cite{bu2025query}。

现有 GraphRAG 工作主要面向大规模文本语料，构建流程计算成本较高，且对文档页面中的图表、表格和视觉布局等多模态信息利用不足。本项目的路线有所不同：直接从页面图像出发，使用 VLM 进行文本/表格/图表/关键词/关系的一体化抽取，构建轻量异构文档图，在查询时综合使用向量检索、图谱邻域扩展和社区补充，是一种更适合课程项目场景的多模态 GraphRAG 实现。

\section{多模态文档理解与问答}

DocVQA 数据集包含 12,000 余张文档图像和 50,000 余个问题，要求模型从文档图像中阅读文字、理解版面结构并给出抽取式答案~\cite{mathew2021docvqa}。ChartQA 聚焦图表问答，涉及柱状图、折线图、趋势比较和简单算术推理~\cite{masry2022chartqa}。这些基准揭示了文档问答的关键特征：视觉布局、表格边界、图表形态和图文对应关系本身就是语义的一部分。

与上述更偏单页图像或特定图表任务的基准相比，pdfQA 更接近本文关注的完整 PDF 问答场景~\cite{schimanski2026pdfqa}。该基准面向真实和合成 PDF 文档，保留了文档级结构、页面顺序、表格、图表、扫描页和复杂版面等信息，并覆盖金融报告、科研论文、杂志、合同和手册等多种文档类型。由于本文系统需要处理多页 ESG 报告中的跨页证据、图表说明和表格信息，因此实验部分选择从 pdfQA 集中抽取报告，并在此基础上构建覆盖文本定位、表格问答、图表理解、跨页延续、多跳关系和引用追踪的自建测试集。

在多模态文档检索方面，ColPali 从页面图像直接生成多向量表示并使用 late interaction 检索，表明页面图像本身可作为检索对象~\cite{faysse2024colpali}。MMDocRAG 基准进一步强调长文档 RAG 需要选择并整合文本、表格、图像等多模态证据~\cite{dong2025mmdocrag}。这些工作启示我们：文档问答系统应保留页面原图、图表描述、bbox 等信息，同时关注证据引用的准确性和完整性。

\section{视觉语言模型}

视觉语言模型（Vision-Language Model, VLM）通过将图像编码器与大语言模型结合，使模型能够在同一语义空间中理解视觉内容和自然语言指令。LLaVA、MiniGPT-4 等工作利用视觉指令微调，使模型具备较强的图像描述、视觉问答和跨模态推理能力~\cite{liu2023llava,zhu2023minigpt4}。在文档理解场景中，VLM 的价值不仅体现在识别页面中的文字，还体现在理解版面结构、表格行列关系、图表趋势、图文对应关系以及跨区域信息组合。相比将 PDF 先拆分为 OCR、版面检测、表格识别和图表解析等多个独立模块，VLM 可以在单次页面理解过程中同时生成文本块、表格、图像说明、实体和关系等结构化结果，从而降低多模块级联带来的误差累积和工程复杂度。

近年来的通用 VLM 也开始针对文档、图表和长上下文任务进行强化。例如 Qwen3-VL 系列支持图像、文本等多模态输入，并在 OCR、空间感知、长上下文理解和多模态推理方面进行了优化~\cite{qwen2025qwen3vl}，适合用于 PDF 页面解析和问答生成。但如果直接将整篇多页 PDF 交给 VLM 进行问答，一方面会带来较高的输入 token 成本，另一方面也容易受到长文档上下文稀释、页间定位困难和证据不可追踪等问题影响。因此，本项目并不将 VLM 作为唯一的端到端问答模块，而是将其分别用于知识库构建阶段的页面结构化解析，以及问答阶段的证据综合与答案生成。

综合来看，纯文本 RAG 结构简单但难以利用图像、表格和版式信息；文本 GraphRAG 能基于实体和关系进行多跳扩展，但仍需要外部模块补充视觉证据；VLM 具备强视觉理解能力，却不适合无筛选地承担整篇长 PDF 的全部检索和推理工作。基于上述分析，本项目采用``FAISS + 异构文档图 + 页面图像 + VLM''的技术路线：先通过向量检索和图结构扩展定位候选证据，再将少量相关页面图像及其结构化节点交给 VLM 生成答案，从而在可追溯性、多模态理解能力和 token 成本之间取得平衡。
